\section[Second order linear difference equations with impulse]{Second order linear dif\/ference equations with impulse}\label{section4}

Consider the second order linear homogeneous dif\/ference equation with impulse%
\begin{gather}
-\Delta ^{2}y_{n-1}+p_{n}y_{n}=0,\qquad n\in \mathbb{Z}_{0}=%
\mathbb{Z}
\backslash \left\{ 0,1\right\} =(-\infty ,-1]\cup \lbrack 2,\infty ),
\label{4.1}
\\
y_{-1}=d_{1}y_{1},\qquad \Delta y_{-1}=d_{2}\Delta y_{1},  \label{4.2}
\end{gather}%
where $y=\left( y_{n}\right) $ with $n\in \mathbb{Z}
$ is a desired solution, the coef\/f\/icients $p_{n}$ are complex numbers given
for $n\in
\mathbb{Z}
_{0}$; $d_{1}$, $d_{2}$ presented in the ``impulse conditions'' (transition
conditions) in (\ref{4.2}) are nonzero complex numbers.

Using the def\/inition of $\Delta $-derivative we can rewrite problem (\ref{4.1}), (\ref{4.2}) in the form%
\begin{gather}
-y_{n-1}+\widetilde{p}_{n}y_{n}-y_{n+1}=0,\qquad n\in (-\infty ,-1]\cup
\lbrack 2,\infty ),  \label{4.3}
\\
y_{-1}=d_{1}y_{1},\qquad y_{0}-y_{-1}=d_{2}(y_{2}-y_{1}),  \label{4.4}
\end{gather}%
where%
\[
\widetilde{p}_{n}=p_{n}+2,\qquad n\in (-\infty ,-1]\cup \lbrack
2,\infty ).
\]

\begin{theorem}
\label{Th4.1}Let $n_{0}$ be a fixed point in $\mathbb{Z}$ and $c_{0}$, $c_{1}$ be given complex numbers. Then problem \eqref{4.1}, \eqref{4.2} has a unique solution $(y_{n})$, $n\in
\mathbb{Z}$, such that%
\begin{equation}
y_{n_{0}}=c_{0},\qquad \Delta y_{n_{0}}=c_{1},\qquad \text{that is,}\qquad
y_{n_{0}}=c_{0},\qquad y_{n_{0}+1}=c_{0}+c_{1}=c_{1}^{\prime }.
\label{4.5}
\end{equation}
\end{theorem}

\begin{proof}
First assume that $n_{0}\in (-\infty ,-1].$ We can rewrite equation (\ref{4.3}) in the form%
\begin{equation}
y_{n-1}=\widetilde{p}_{n}y_{n}-y_{n+1}=0,\qquad n\in (-\infty ,-1]\cup
\lbrack 2,\infty )  \label{4.6}
\end{equation}%
as well as in the form%
\begin{equation}
y_{n+1}=\widetilde{p}_{n}y_{n}-y_{n-1}=0,\qquad n\in (-\infty ,-1]\cup
\lbrack 2,\infty ).  \label{4.7}
\end{equation}%
Using the initial conditions (\ref{4.5}) we f\/ind, recurrently (step by
step), $y_{n}$ for $n\leq n_{0}+1$ uniquely from (\ref{4.6}) and for $%
n_{0}+2\leq n\leq -1$ uniquely from (\ref{4.7}). Then we f\/ind $y_{1}$ and $%
y_{2}$ from the transition conditions (\ref{4.4}) and then we f\/ind $y_{n}$
for $n\geq 3$ uniquely from (\ref{4.7}).

In the case $n_{0}\in \lbrack 1,\infty )$ we are reasoning similarly; using
equations (\ref{4.6}), (\ref{4.7}) we f\/irst f\/ind~$y_{n}$ uniquely for $n\geq
1$ and then using the transition conditions (\ref{4.4}) we pass to the
interval $(-\infty ,-1].$

Finally, if $n_{0}=0,$ then we f\/ind $y_{0}$ and $y_{1}$ uniquely from the
initial conditions (\ref{4.5}) with $n_{0}=0$. Then we f\/ind $y_{-1}$ and $%
y_{2}$ from the transition conditions (\ref{4.4}). Next, solving equation~(\ref{4.6}) at f\/irst on $(-\infty ,-1]$ we f\/ind $y_{n}$ uniquely for $n\in
(-\infty ,-2]$ and then solving (\ref{4.7}) on $[2,\infty )$ we f\/ind $y_{n}$
uniquely for $n\in \lbrack 3,\infty ).$
\end{proof}

\begin{definition}
\label{Def4.2}For two sequences $y=(y_{n})$ and $z=(z_{n})$ with $n\in
\mathbb{Z}
$, we def\/ine their Wronskian~by%
\begin{gather*}
W_{n}(y,z) = y_{n}\Delta z_{n}-(\Delta y_{n})z_{n}  = y_{n}z_{n+1}-y_{n+1}z_{n},\qquad n\in
\mathbb{Z}
.
\end{gather*}
\end{definition}

\begin{theorem}
\label{Th4.3} The Wronskian of any two solutions $y$ and $z$ of problem \eqref{4.1}, \eqref{4.2} is constant on each of the intervals $(-\infty ,-1]$ and
$[1,\infty )$:
\begin{equation}
W_{n}(y,z)=\left\{
\begin{array}{lll}
\omega ^{-} & \text{if} & n\in (-\infty ,-1], \\
\omega ^{+} & \text{if} & n\in \lbrack 1,\infty ).%
\end{array}%
\right.  \label{4.8}
\end{equation}%
In addition,
\begin{equation}
\omega ^{-}=d_{1}d_{2}\omega ^{+}  \label{4.9}
\end{equation}%
and%
\begin{equation}
W_{0}(y,z)=-d_{2}\omega ^{+}.  \label{4.10}
\end{equation}
\end{theorem}

\begin{proof}
Suppose that $y=(y_{n})$ and $z=(z_{n}),$ where $n\in
\mathbb{Z}
,$ are solutions of (\ref{4.1}), (\ref{4.2}). Let us compute the $\Delta $%
-derivative of $W_{n}(y,z)$. Using the product rule for $\Delta $-derivative%
\begin{gather*}
\Delta (f_{n}g_{n}) = (\Delta f_{n})g_{n}+f_{n+1}\Delta g_{n}
= f_{n}\Delta g_{n}+(\Delta f_{n})g_{n+1},
\end{gather*}
we have
\begin{gather*}
\Delta W_{n}(y,z)=\Delta \left[ y_{n}\Delta z_{n}-(\Delta y_{n})z_{n}\right]
=(\Delta y_{n})\Delta z_{n}+y_{n+1}\Delta ^{2}z_{n}-(\Delta y_{n})\Delta
z_{n}-(\Delta ^{2}y_{n})z_{n+1}
\\
\phantom{\Delta W_{n}(y,z)}{}
=y_{n+1}\Delta ^{2}z_{n}-(\Delta ^{2}y_{n})z_{n+1}.
\end{gather*}
Further, since $y_{n}$ and $z_{n}$ are solutions of (\ref{4.1}), (\ref{4.2}),%
\begin{gather*}
\Delta ^{2}y_{n}=p_{n+1}y_{n+1},\qquad n\in (-\infty ,-2]\cup \lbrack
1,\infty ),
\\
\Delta ^{2}z_{n}=p_{n+1}z_{n+1},\qquad n\in (-\infty ,-2]\cup \lbrack
1,\infty ).
\end{gather*}
Therefore%
\[
\Delta W_{n}(y,z)=0\qquad \text{for} \ \ n\in (-\infty ,-2]\cup \lbrack
1,\infty ).
\]%
The latter implies that $W_{n}(y,z)$ is constant on $(-\infty ,-1]$ and on $%
[1,\infty ).$ Thus we have (\ref{4.8}), where $\omega ^{-}$ and $\omega ^{+}$
are some constants (depending on the solutions $y$ and $z$).

Next using (\ref{4.8}) and the impulse conditions in (\ref{4.2}) for $y_{n}$
and $z_{n},$ we have{\samepage
\begin{gather*}
\omega ^{-}=W_{-1}(y,z)=y_{-1}\Delta z_{-1}-(\Delta y_{-1})z_{-1}
=d_{1}d_{2}[y_{1}\Delta z_{1}-(\Delta
y_{1})z_{1}]\\
\phantom{\omega ^{-}}{} =d_{1}d_{2}W_{1}(y,z)=d_{1}d_{2}\omega ^{+},
\end{gather*}
so that (\ref{4.9}) is established.}

Finally, from the impulse conditions in (\ref{4.2}) we f\/ind that%
\[
y_{0}=(d_{1}-d_{2})y_{1}+d_{2}y_{2}.
\]%
Substituting this expression for $y_{0}$ and $z_{0}$ into
\[
W_{0}(y,z)=y_{0}z_{1}-y_{1}z_{0},
\]%
we get%
\[
W_{0}(y,z)=-d_{2}W_{1}(y,z)=-d_{2}\omega ^{+}.
\]%
Therefore (\ref{4.10}) is also proved.
\end{proof}

\begin{corollary}
\label{Cor4.4}If $y$ and $z$ are two solutions of \eqref{4.1}, \eqref{4.2},
then either $W_{n}(y,z)=0$ for all $n\in
\mathbb{Z}$ or $W_{n}(y,z)\neq 0$ for all $n\in
\mathbb{Z}.$
\end{corollary}

By using Theorem \ref{Th4.1}, the following two theorems can be proved in
exactly the same way when equation (\ref{4.1}) does not include any impulse
conditions \cite{[6]}.

\begin{theorem}
\label{Th4.5}Any two solutions of \eqref{4.1}, \eqref{4.2} are linearly
independent if and only if their Wronskian is not zero.
\end{theorem}

\begin{theorem}
\label{Th4.6}Problem \eqref{4.1}, \eqref{4.2} has two linearly independent
solutions and every solution of~\eqref{4.1},~\eqref{4.2} is a linear
combination of these solutions.
\end{theorem}

We say that $y=(y_{n})$ and $z=(z_{n}),$ where $n\in
%TCIMACRO{\U{2124} }%
%BeginExpansion
\mathbb{Z}
%EndExpansion
,$ form a \textit{fundamental set }(or \textit{fundamental system}) of
solutions for (\ref{4.1}), (\ref{4.2}) provided that they are solutions of (%
\ref{4.1}), (\ref{4.2}) and their Wronskian is not zero.

Let us consider the nonhomogeneous equation%
\begin{equation}
-\Delta ^{2}y_{n-1}+p_{n}y_{n}=h_{n},\qquad n\in (-\infty ,-1]\cup
\lbrack 2,\infty ),  \label{4.11}
\end{equation}%
with the impulse conditions%
\begin{equation}
y_{-1}=d_{1}y_{1},\qquad \Delta y_{-1}=d_{2}\Delta y_{1},  \label{4.12}
\end{equation}%
where $h_{n}$ is a complex sequence def\/ined for $n\in (-\infty ,-1]\cup
\lbrack 2,\infty ).$ We will extend $h_{n}$ to the values $n=0$ and $n=1$ by
setting%
\begin{equation}
h_{0}=h_{1}=0.  \label{4.13}
\end{equation}

\begin{theorem}
\label{Th4.7}Suppose that $u=(u_{n})$ and $v=(v_{n})$ form a fundamental set
of solutions of the homogeneous problem \eqref{4.1}, \eqref{4.2}. Then a
general solution of the corresponding nonhomogeneous problem \eqref{4.11}, \eqref{4.12} is given by%
\[
y_{n}=c_{1}u_{n}+c_{2}v_{n}+x_{n},\qquad n\in
\mathbb{Z},
\]%
where $c_{1}$, $c_{2}$ are arbitrary constants and%
\begin{equation}
x_{n}=\left\{
\begin{array}{lll}
\displaystyle -\sum_{s=n}^{0}\frac{u_{n}v_{s}-u_{s}v_{n}}{W_{s}(u,v)}h_{s} & \text{if} &
n\leq 0, \vspace{1mm} \\
\displaystyle \sum_{s=1}^{n}\frac{u_{n}v_{s}-u_{s}v_{n}}{W_{s}(u,v)}h_{s} & \text{if} &
n\geq 1.%
\end{array}%
\right.  \label{4.14}
\end{equation}
\end{theorem}

\begin{proof}
Taking into account (\ref{4.13}) it is not dif\/f\/icult to verify that the
sequence $x_{n}$ def\/ined by~(\ref{4.14}) is a particular solution of (\ref%
{4.11}), (\ref{4.12}), namely, $x_{n}$ satisf\/ies equation (\ref{4.11}) and
the conditions%
\[
x_{-1}=\Delta x_{-1}=0,\qquad x_{1}=\Delta x_{1}=0.
\]%
This implies that the statement of the theorem is true.
\end{proof}

\section{Two special solutions}\label{section5}

Consider the homogeneous problem%
\begin{gather}
-\Delta ^{2}y_{n-1}+q_{n}y_{n}=0,\qquad n\in
\mathbb{Z}
_{0}=
\mathbb{Z}
\backslash \{0,1\},  \label{5.1}
\\
y_{-1}=y_{1},\qquad \Delta y_{-1}=e^{2i\delta }\Delta y_{1},
\label{5.2}
\end{gather}%
where $\delta \in \lbrack 0,\pi /2)$ is a f\/ixed real number and%
\begin{equation}
q_{n}\geq c>0\qquad \text{for} \qquad n\in
\mathbb{Z}
_{0}.  \label{5.3}
\end{equation}

In this section we show that under the condition (\ref{5.3}) problem (\ref%
{5.1}), (\ref{5.2}) has two linearly independent solutions $\psi =(\psi
_{n}) $ and $\chi =(\chi _{n}),$ where $n\in
\mathbb{Z}
,$ such that%
\begin{equation}
\sum_{n=0}^{\infty }\left\vert \psi _{n}\right\vert ^{2}<\infty \qquad \text{and} \qquad \sum_{-\infty }^{n=0}\left\vert \chi _{n}\right\vert ^{2}<\infty .
\label{5.4}
\end{equation}%
These solutions will allow us to f\/ind the inverse $L^{-1}$ of the operator $%
L $ introduced above in Section~\ref{section2} and investigate the properties of
$L^{-1}.$

First we derive two simple useful formulas related to the nonhomogeneous
problem%
\begin{gather}
-\Delta ^{2}y_{n-1}+q_{n}y_{n}=f_{n},\qquad n\in
\mathbb{Z}
_{0},  \label{5.5}
\\
y_{-1}=y_{1},\qquad \Delta y_{-1}=e^{2i\delta }\Delta y_{1},
\label{5.6}
\end{gather}
where $(q_{n})$ is a real sequence with $n\in
\mathbb{Z}
_{0},$ and $\delta \in \lbrack 0,\pi /2);$ $(f_{n})$ is a complex sequence
with $n\in
\mathbb{Z}
_{0}.$

\begin{lemma}
\label{Lem5.1}Let $y=(y_{n})$ with $n\in
\mathbb{Z}
$ be a solution of problem \eqref{5.5}, \eqref{5.6} and $a,$ $b$ be any
integers such that $a\leq -1$ and $b\geq 2$. Then the following formulas
hold:
\begin{gather}
\sum_{n=2}^{b}\big( \left\vert \Delta y_{n}\right\vert ^{2}+q_{n}\left\vert
y_{n}\right\vert ^{2}\big) =(\Delta y_{n})\overline{y}_{n+1}\big|
_{1}^{b}+\sum_{n=2}^{b}f_{n}\overline{y}_{n},  \label{5.7}
\\
\sum_{n=a}^{-1}\big( \left\vert \Delta y_{n}\right\vert
^{2}+q_{n}\left\vert y_{n}\right\vert ^{2}\big) =(\Delta y_{n})\overline{y}%
_{n+1}\big|_{a-1}^{-1}+\sum_{n=a}^{-1}f_{n}\overline{y}_{n}.  \label{5.8}
\end{gather}
\end{lemma}

\begin{proof}
To prove (\ref{5.7}), multiply equation (\ref{5.5}) by $\overline{y}_{n}$
and sum from $n=2$ to $n=b$:
\[
-\sum_{n=2}^{b}(\Delta ^{2}y_{n-1})\overline{y}_{n}+\sum_{n=2}^{b}q_{n}\left%
\vert y_{n}\right\vert ^{2}=\sum_{n=2}^{b}f_{n}\overline{y}_{n}.
\]%
Next, applying the summation by parts formula (\ref{3.1}) we get that%
\[
-\sum_{n=2}^{b}(\Delta ^{2}y_{n-1})\overline{y}_{n}=-\sum_{n=2}^{b}(\nabla
\Delta y_{n})\overline{y}_{n}=-(\Delta y_{n})\overline{y}_{n+1}\big|
_{1}^{b}+\sum_{n=2}^{b}\left\vert \Delta y_{n}\right\vert ^{2}.
\]%
Therefore the formula (\ref{5.7}) follows.

The formula (\ref{5.8}) can be proved in a similar way.
\end{proof}

\begin{theorem}
\label{Th5.2}Under the condition \eqref{5.3} problem \eqref{5.1}, \eqref{5.2} has two linearly independent solutions $\psi =(\psi _{n})$ and $\chi =(\chi
_{n})$ with $n\in
\mathbb{Z}
$, possessing the properties stated in \eqref{5.4}.
\end{theorem}

\begin{proof}
Denote by $\varphi =(\varphi _{n})$ and $\theta =(\theta _{n}),$ where $n\in
\mathbb{Z}
,$ solutions of problem (\ref{5.1}), (\ref{5.2}) satisfying the initial
conditions%
\begin{gather}
\varphi _{1}=1,\qquad \Delta \varphi _{1}=-1,  \label{5.9}
\\
\theta _{1}=1,\qquad \Delta \theta _{1}=0.  \label{5.10}
\end{gather}%
Such solutions exist and are unique by Theorem~\ref{Th4.1}. It follows from (%
\ref{5.9}), (\ref{5.10}) that $\varphi _{2}=0,$ $\theta _{2}=1.$ According
to Theorem \ref{Th4.3} we f\/ind that%
\begin{equation}
W_{0}(\varphi ,\theta )=-e^{2i\delta },\qquad W_{n}(\varphi ,\theta
)=\left\{
\begin{array}{lll}
e^{2i\delta } & \text{if} & n\leq -1, \\
1 & \text{if} & n\geq 1.%
\end{array}%
\right.  \label{5.11}
\end{equation}%
Therefore $W_{n}(\varphi ,\theta )\neq 0$ and by Theorem \ref{Th4.5} the
solutions $\varphi $ and $\theta $ are linearly independent.

We seek the desired solution $\psi =(\psi _{n})$ of problem (\ref{5.1}), (\ref{5.2}) in the form%
\begin{equation}
\psi _{n}=\varphi _{n}+v\theta _{n},\qquad n\in
\mathbb{Z}
,  \label{5.12}
\end{equation}%
where $v$ is a complex constant which we will choose.

Take an arbitrary integer $b\geq 2$. Applying (\ref{5.7}) to%
\begin{gather*}
-\Delta ^{2}\psi _{n-1}+q_{n}\psi _{n}=0,\qquad n\in
\mathbb{Z}
_{0},
\\
\psi _{-1}=\psi _{1},\qquad \Delta \psi _{-1}=e^{2i\delta }\Delta \psi
_{1},
\end{gather*}
we get%
\begin{equation*}
\sum_{n=2}^{b}\big( \left\vert \Delta \psi _{n}\right\vert
^{2}+q_{n}\left\vert \psi _{n}\right\vert ^{2}\big) =(\Delta \psi _{n})%
\overline{\psi }_{n+1}\big|_{1}^{b}.  %\label{5.13}
\end{equation*}%
Since $\Delta \psi _{1}=-1$ and $\psi _{2}=v$, by (\ref{5.12}) and (\ref{5.9}), (\ref{5.10}), hence%
\[
\sum_{n=2}^{b}\big( \left\vert \Delta \psi _{n}\right\vert
^{2}+q_{n}\left\vert \psi _{n}\right\vert ^{2}\big) =(\Delta \psi _{b})%
\overline{\psi }_{b+1}+\overline{v}.
\]%
Multiply the latter equality by $e^{i\delta }$ and take then the real part
of both sides to get%
\begin{equation}
(\cos \delta )\sum_{n=2}^{b}\big( \left\vert \Delta \psi _{n}\right\vert
^{2}+q_{n}\left\vert \psi _{n}\right\vert ^{2}\big) ={\rm Re}\big\{e^{i\delta
}(\Delta \psi _{b})\overline{\psi }_{b+1}\big\}+{\rm Re}\big(ve^{-i\delta }\big).
\label{5.14}
\end{equation}

Now we choose $v$ so that to have%
\begin{equation}
{\rm Re}\big\{e^{i\delta }(\Delta \psi _{b})\overline{\psi }_{b+1}\big\}=0.
\label{5.15}
\end{equation}%
Since%
\[
{\rm Re}\big\{e^{i\delta }(\Delta \psi _{b})\overline{\psi }_{b+1}\big\}=\left\vert
\psi _{b+1}\right\vert ^{2}{\rm Re}\left\{ e^{i\delta }\frac{\Delta \psi
_{b}}{\psi _{b+1}}\right\} ,
\]%
it is suf\/f\/icient for (\ref{5.15}) to have%
\begin{equation}
{\rm Re}\left\{ e^{i\delta }\frac{\Delta \psi _{b}}{\psi _{b+1}}\right\} =0.
\label{5.16}
\end{equation}%
Note that $\psi _{b}$ cannot be zero for any two successive values of $b$
(otherwise $\psi _{n}$ would be identically zero by the uniqueness theorem
for solution, that is not true since $\varphi $ and $\theta $ are linearly
independent). Therefore $\psi _{b}\neq 0$ for inf\/initely many values of $b.$

Under the condition (\ref{5.16}) the equation (\ref{5.14}) becomes%
\begin{equation*}
(\cos \delta )\sum_{n=2}^{b}\big( \left\vert \Delta \psi _{n}\right\vert
^{2}+q_{n}\left\vert \psi _{n}\right\vert ^{2}\big) ={\rm Re}
\big(ve^{-i\delta }\big).  %\label{5.17}
\end{equation*}%
The condition (\ref{5.16}) can be written as%
\begin{equation}
e^{i\delta }\frac{\Delta \varphi _{b}+v\Delta \theta _{b}}{\varphi
_{b+1}+v\theta _{b+1}}=\beta ,  \label{5.18}
\end{equation}%
where $\beta $ is a pure imaginary number ($\beta =it,$ $t\in
\mathbb{R}
)$. Note that
\begin{equation}
\varphi _{b+1}\Delta \theta _{b}-(\Delta \varphi _{b})\theta _{b+1}=-\varphi
_{b+1}\theta _{b}+\varphi _{b}\theta _{b+1}=W_{b}(\varphi ,\theta )=1\neq 0
\label{5.19}
\end{equation}
by (\ref{5.11}). Therefore (\ref{5.18}) def\/ines a linear-fractional
transformation of the complex $v$-plane onto the complex $\beta $-plane.
Solving (\ref{5.18}) for $v,$ we get%
\begin{equation}
v(\beta )=\frac{\varphi _{b+1}\beta -e^{i\delta }\Delta \varphi _{b}}{%
-\theta _{b+1}\beta +e^{i\delta }\Delta \theta _{b}}.  \label{5.20}
\end{equation}%
Thus, condition (\ref{5.16}) will be satisf\/ied if we choose $v$ by (\ref%
{5.20}) for pure imaginary values of $\beta .$ On the other hand, when $%
\beta $ runs in (\ref{5.20}) the imaginary axis, $v(\beta )$ describes a
circle $C_{b}$ in the $v$-plane. The center of the circle is symmetric point
of the point at inf\/inity with respect to the circle. Since%
\[
v(\beta ^{\prime })=\infty ,\qquad \text{where} \qquad \beta ^{\prime
}=e^{i\delta }\frac{\Delta \theta _{b}}{\theta _{b+1}},
\]%
the point%
\[
\beta _{0}=-\overline{\beta ^{\prime }}=-e^{-i\delta }\frac{\Delta \overline{%
\theta }_{b}}{\overline{\theta }_{b+1}}
\]%
which is symmetric point of the point $\beta ^{\prime }$ with respect to the
imaginary axis of the $\beta $-plane, is mapped onto the center of $C_{b}$.
So the center of $C_{b}$ is located at the point%
\begin{equation}
v(\beta _{0})=-\frac{e^{-i\delta }\varphi _{b+1}\Delta \overline{\theta }%
_{b}+e^{i\delta }(\Delta \varphi _{b})\overline{\theta }_{b+1}}{e^{-i\delta
}\theta _{b+1}\Delta \overline{\theta }_{b}+e^{i\delta }\overline{\theta }%
_{b+1}\Delta \theta _{b}}.  \label{5.21}
\end{equation}%
Note that the denominator in (\ref{5.21}) is dif\/ferent from zero. This fact
follows from the equality%
\begin{gather}
e^{-i\delta }\theta _{b+1}\Delta \overline{\theta }_{b}+e^{i\delta }%
\overline{\theta }_{b+1}\Delta \theta _{b}=2{\rm Re}\big\{ e^{i\delta
}(\Delta \theta _{n})\overline{\theta }_{n+1}\big|_{1}^{b}\big\}\nonumber\\
\phantom{e^{-i\delta }\theta _{b+1}\Delta \overline{\theta }_{b}+e^{i\delta }%
\overline{\theta }_{b+1}\Delta \theta _{b}}{}
=(2\cos \delta )\sum_{n=2}^{b}\big( \left\vert \Delta \theta
_{n}\right\vert ^{2}+q_{n}\left\vert \theta _{n}\right\vert ^{2}\big),
\label{5.22}
\end{gather}%
which can be derived as (\ref{5.14}) taking into account (\ref{5.10}). The
radius $R_{b}$ of the circle $C_{b}$ is equal to the distance between the
center $v(\beta _{0})$ of $C_{b}$ and the point $v(0)$ on the circle.
Calculating the dif\/ference $v(\beta _{0})-v(0)$ by using (\ref{5.19})--(\ref{5.22}) we easily f\/ind that%
\[
R_{b}=\frac{1}{(2\cos \delta )\sum\limits_{n=2}^{b}\big( \left\vert \Delta \theta
_{n}\right\vert ^{2}+q_{n}\left\vert \theta _{n}\right\vert ^{2}\big) }.
\]%
Further, since%
\[
{\rm Re}(e^{i\delta }\overline{\theta }_{b+1}\Delta \theta
_{b})=-\left\vert \theta _{b+1}\right\vert ^{2}{\rm Re}\,\beta _{0},
\]%
we get from (\ref{5.22})%
\[
(\cos \delta )\sum_{n=2}^{b}\big( \left\vert \Delta \theta _{n}\right\vert
^{2}+q_{n}\left\vert \theta _{n}\right\vert ^{2}\big) =-\left\vert \theta
_{b+1}\right\vert ^{2}{\rm Re}\, \beta _{0}.
\]%
Therefore ${\rm Re}\,\beta _{0}<0.$ This means that the left half-plane of
the $\beta $-plane is mapped onto the interior of the circle $C_{b}.$
Consequently, $v(\beta )$ belongs to the interior of the circle $C_{b}$ if
and only if ${\rm Re}\,\beta <0.$ This inequality is equivalent by (\ref{5.14}), (\ref{5.18}) to%
\begin{equation}
(\cos \delta )\sum_{n=2}^{b}\big( \left\vert \Delta \psi _{n}\right\vert
^{2}+q_{n}\left\vert \psi _{n}\right\vert ^{2}\big) <{\rm Re}\big(ve^{-i\delta }\big).  \label{5.23}
\end{equation}

Thus, $v$ belongs to the interior of the circle $C_{b}$ if and only if the
inequality (\ref{5.23}) holds and~$v$ lies on the circle $C_{b}$ if and only
if%
\[
(\cos \delta )\sum_{n=2}^{b}\big( \left\vert \Delta \psi _{n}\right\vert
^{2}+q_{n}\left\vert \psi _{n}\right\vert ^{2}\big) ={\rm Re}
\big(ve^{-i\delta }\big).
\]

Now let $b_{2}>b_{1}.$ Then if $v$ is inside or on $C_{b_{2}}$%
\[
(\cos \delta )\sum_{n=2}^{b_{1}}\big( \left\vert \Delta \psi
_{n}\right\vert ^{2}+q_{n}\left\vert \psi _{n}\right\vert ^{2}\big) <(\cos
\delta )\sum_{n=2}^{b_{2}}\big( \left\vert \Delta \psi _{n}\right\vert
^{2}+q_{n}\left\vert \psi _{n}\right\vert ^{2}\big) \leq {\rm Re}
\big(ve^{-i\delta }\big)
\]%
and therefore $v$ is inside $C_{b_{1}}.$ This means $C_{b_{1}}$ contains $%
C_{b_{2}}$ in its interior if $b_{2}>b_{1}.$ It follows that, as $%
b\rightarrow \infty ,$ the circles $C_{b}$ converge either to a limit-circle
or to a limit-point. If $\widehat{v}$ is the limit-point or any point on the
limit-circle, then $\widehat{v}$ is inside any $C_{b}.$ Hence%
\[
(\cos \delta )\sum_{n=2}^{b}\big( \left\vert \Delta \psi _{n}\right\vert
^{2}+q_{n}\left\vert \psi _{n}\right\vert ^{2}\big) <{\rm Re}\big(\widehat{v}%
e^{-i\delta }\big),
\]%
where%
\begin{equation}
\psi _{n}=\varphi _{n}+\widehat{v}\theta _{n},\qquad n\in
\mathbb{Z}
,  \label{5.24}
\end{equation}%
and letting $b\rightarrow \infty $ we get%
\begin{equation}
(\cos \delta )\sum_{n=2}^{\infty }\big( \left\vert \Delta \psi
_{n}\right\vert ^{2}+q_{n}\left\vert \psi _{n}\right\vert ^{2}\big) \leq
{\rm Re}\big(\widehat{v}e^{-i\delta }\big).  \label{5.25}
\end{equation}%
It also follows that%
\begin{equation}
{\rm Re}\big(\widehat{v}e^{-i\delta }\big)>0.  \label{5.26}
\end{equation}%
Since $q_{n}\geq c>0,$ (\ref{5.25}) implies that for the solution $\psi
=(\psi _{n})$ def\/ined by (\ref{5.24}) we have (\ref{5.4}). Thus, the
statement of the theorem concerning the solution $\psi =(\psi _{n})$ is
proved.

Let us now show existence of the solution $\chi =(\chi _{n})$ satisfying (%
\ref{5.4}). We seek the desired solution $\theta =(\theta _{n})$ of problem (%
\ref{5.1}), (\ref{5.2}) in the form%
\[
\chi _{n}=\varphi _{n}+u\theta _{n},\qquad n\in
\mathbb{Z}
,
\]
where $u$ is a complex constant to be determined.

Take an arbitrary integer $a\leq -1.$ Applying (\ref{5.8}) to the equations%
\begin{gather*}
-\Delta ^{2}\chi _{n-1}+q_{n}\chi _{n}=0,\qquad n\in
\mathbb{Z}
_{0},
\\
\chi _{-1}=\chi _{1},\qquad \Delta \chi _{-1}=e^{2i\delta }\Delta \chi
_{1},
\end{gather*}
we get%
\[
\sum_{n=a}^{-1}\big( \left\vert \Delta \chi _{n}\right\vert
^{2}+q_{n}\left\vert \chi _{n}\right\vert ^{2}\big) =(\Delta \chi _{n})%
\overline{\chi }_{n+1}\big|_{a-1}^{-1}.
\]%
Since%
\[
\Delta \chi _{-1}=-e^{2i\delta },\qquad \chi _{0}=1-e^{2i\delta }+u,
\]%
we have%
\begin{equation}
\sum_{n=a}^{-1}\big( \left\vert \Delta \chi _{n}\right\vert
^{2}+q_{n}\left\vert \chi _{n}\right\vert ^{2}\big) =-e^{2i\delta }+1-%
\overline{u}e^{2i\delta }-(\Delta \chi _{a-1})\overline{\chi }_{a}.
\label{5.27}
\end{equation}%
Multiply both sides of (\ref{5.27}) by $e^{-i\delta }$ and take then the
real part of both sides to get%
\begin{equation}
(\cos \delta )\sum_{n=a}^{-1}\big( \left\vert \Delta \chi _{n}\right\vert
^{2}+q_{n}\left\vert \chi _{n}\right\vert ^{2}\big) =-{\rm Re}
\big(ue^{-i\delta }\big)-{\rm Re}\big\{ e^{-i\delta }(\Delta \chi _{a-1})%
\overline{\chi }_{a}\big\} .  \label{5.28}
\end{equation}

Now we choose $u$ so that to have%
\begin{equation}
{\rm Re}\big\{e^{-i\delta }(\Delta \chi _{a-1})\overline{\chi }_{a}\big\}=0.
\label{5.29}
\end{equation}%
Since%
\[
{\rm Re}\big\{e^{-i\delta }(\Delta \chi _{a-1})\overline{\chi }%
_{a}\big\}=\left\vert \chi _{a}\right\vert ^{2}{\rm Re}\left\{ e^{-i\delta }%
\frac{\Delta \chi _{a-1}}{\chi _{a}}\right\} ,
\]%
it is suf\/f\/icient for (\ref{5.29}) to have%
\begin{equation}
{\rm Re} \left\{ e^{-i\delta }\frac{\Delta \chi _{a-1}}{\chi _{a}}\right\}
=0.  \label{5.30}
\end{equation}%
Then (\ref{5.28}) becomes
\begin{equation*}
(\cos \delta )\sum_{n=a}^{-1}\left( \left\vert \Delta \chi _{n}\right\vert
^{2}+q_{n}\left\vert \chi _{n}\right\vert ^{2}\right) =-{\rm Re}\big(ue^{-i\delta }\big).  %\label{5.31}
\end{equation*}%
The condition (\ref{5.30}) can be written as%
\begin{equation}
e^{-i\delta }\frac{\Delta \varphi _{a-1}+u\Delta \theta _{a-1}}{\varphi
_{a}+u\theta _{a}}=\alpha ,  \label{5.32}
\end{equation}%
where $\alpha $ is a pure imaginary number ($\alpha =it$, $t\in
\mathbb{R}).$ Note that%
\begin{equation}
\varphi _{a}\Delta \theta _{a-1}-(\Delta \varphi _{a-1})\theta
_{a}=W_{a-1}(\varphi ,\theta )=e^{2i\delta }\neq 0  \label{5.33}
\end{equation}%
by (\ref{5.11}). Therefore (\ref{5.32}) def\/ines a linear-fractional
transformation of the complex $u$-plane onto the complex $\alpha $-plane.
Solving (\ref{5.32}) for $u,$ we get%
\begin{equation}
u(\alpha )=\frac{\varphi _{a}\alpha -e^{-i\delta }\Delta \varphi _{a-1}}{%
-\theta _{a}\alpha +e^{-i\delta }\Delta \theta _{a-1}}.  \label{5.34}
\end{equation}%
Thus, condition (\ref{5.30}) will be satisf\/ied if we choose $u$ by (\ref%
{5.34}) for pure imaginary values of $\alpha .$ On the other hand, when $%
\alpha $ runs in (\ref{5.34}) the imaginary axis, $u(\alpha )$ describes a
circle $K_{a}$ in the $u$-plane. The center of the circle is symmetric point
of the point at inf\/inity with respect to the circle. Since%
\[
u(\alpha ^{\prime })=\infty ,\qquad \text{where} \qquad \alpha ^{\prime
}=e^{-i\delta }\frac{\Delta \theta _{a-1}}{\theta _{a}},
\]%
the point%
\[
\alpha _{0}=-\overline{\alpha ^{\prime }}=-e^{i\delta }\frac{\Delta
\overline{\theta }_{a-1}}{\overline{\theta }_{a}}
\]%
which is symmetric point of the point $\alpha ^{\prime }$ with respect to
the imaginary axis of the $\alpha $-plane, is mapped onto the center of $%
K_{a}.$ So the center of $K_{a}$ is located at the point%
\begin{equation}
u(\alpha _{0})=-\frac{e^{i\delta }\varphi _{a}\Delta \overline{\theta }%
_{a-1}+e^{-i\delta }(\Delta \varphi _{a-1})\overline{\theta }_{a}}{%
e^{i\delta }\theta _{a}\Delta \overline{\theta }_{a-1}+e^{-i\delta }%
\overline{\theta }_{a}\Delta \theta _{a-1}}.  \label{5.35}
\end{equation}%
Note that the denominator in (\ref{5.35}) is dif\/ferent from zero. This fact
follows from the equality%
\begin{gather}
e^{-i\delta }\theta _{a}\Delta \overline{\theta }_{a-1}+e^{i\delta }%
\overline{\theta }_{a}\Delta \theta _{a-1}=2{\rm Re}\big\{ e^{i\delta
}(\Delta \theta _{a-1})\overline{\theta }_{a}\big\}\nonumber\\
\qquad {}=-(2\cos \delta )\sum_{n=a}^{-1}\big( \left\vert \Delta \theta
_{n}\right\vert ^{2}+q_{n}\left\vert \theta _{n}\right\vert ^{2}\big)
\label{5.36}
\end{gather}%
which can be derived as (\ref{5.28}) taking into account $\Delta \theta
_{-1}=0,$ $\theta _{0}=1.$ Calculating the dif\/ference $u(\alpha _{0})-u(0)$
we easily f\/ind the radius $R_{a}=\left\vert u(\alpha _{0})-u(0)\right\vert $
of the circle $K_{a},$ using (\ref{5.33})--(\ref{5.36}),%
\[
R_{a}=\frac{1}{(2\cos \delta )\sum\limits_{n=a}^{-1}\big( \left\vert \Delta \theta
_{n}\right\vert ^{2}+q_{n}\left\vert \theta _{n}\right\vert ^{2}\big) }.
\]%
Further, since%
\[
{\rm Re}\big(e^{-i\delta }\overline{\theta }_{a}\Delta \theta
_{a-1}\big)=-\left\vert \theta _{a}\right\vert ^{2}{\rm Re}\,\alpha _{0},
\]%
we get from (\ref{5.36})%
\[
(\cos \delta )\sum_{n=a}^{-1}\big( \left\vert \Delta \theta _{n}\right\vert
^{2}+q_{n}\left\vert \theta _{n}\right\vert ^{2}\big) =\left\vert \theta
_{a}\right\vert ^{2}{\rm Re}\,\alpha _{0}.
\]%
Therefore ${\rm Re}\,\alpha _{0}>0.$ This means that the right half-plane
of the $\alpha $-plane is mapped onto the interior of the circle $K_{a}.$
Consequently, $u(\alpha )$ lies inside the circle $K_{a}$ if and only if $%
{\rm Re}\,\alpha >0.$ This inequality is equivalent by (\ref{5.28}), (\ref{5.32}) to
\begin{equation}
(\cos \delta )\sum_{n=a}^{-1}\big( \left\vert \Delta \chi _{n}\right\vert
^{2}+q_{n}\left\vert \chi _{n}\right\vert ^{2}\big) <-{\rm Re}
\big(ue^{-i\delta }\big).  \label{5.37}
\end{equation}

Thus, $u$ lies inside the circle $K_{a}$ if and only if the inequality (\ref%
{5.37}) holds and $u$ lies on the circle $K_{a}$ if and only if%
\[
(\cos \delta )\sum_{n=a}^{-1}\big( \left\vert \Delta \chi _{n}\right\vert
^{2}+q_{n}\left\vert \chi _{n}\right\vert ^{2}\big) =-{\rm Re}
\big(ue^{-i\delta }\big).
\]

Now let $a_{2}<a_{1}.$ Then if $u$ is inside or on $K_{a_{2}}$%
\[
(\cos \delta )\sum_{n=a_{1}}^{-1}\big( \left\vert \Delta \chi
_{n}\right\vert ^{2}+q_{n}\left\vert \chi _{n}\right\vert ^{2}\big) <(\cos
\delta )\sum_{n=a_{2}}^{-1}\big( \left\vert \Delta \chi _{n}\right\vert
^{2}+q_{n}\left\vert \chi _{n}\right\vert ^{2}\big) \leq -{\rm Re}
\big(ue^{-i\delta }\big)
\]%
and therefore $u$ is inside $K_{a_{1}}$. This means $K_{a_{1}}$ contains $%
K_{a_{2}}$ in its interior if $a_{2}<a_{1}.$ It follows that, as $%
a\rightarrow -\infty ,$ the circles $K_{a}$ converge either to a
limit-circle or to a limit-point. If $\widehat{u}$ is the limit-point or any
point on the limit-circle, then $\widehat{u}$ is inside any $K_{a}.$ Hence%
\[
(\cos \delta )\sum_{n=a}^{-1}\big( \left\vert \Delta \chi _{n}\right\vert
^{2}+q_{n}\left\vert \chi _{n}\right\vert ^{2}\big) <-{\rm Re}\big(\widehat{u}%
e^{-i\delta }\big),
\]%
where%
\begin{equation}
\chi _{n}=\varphi _{n}+\widehat{u}\theta _{n},\qquad n\in
\mathbb{Z}
,  \label{5.38}
\end{equation}%
and letting $a\rightarrow -\infty $ we get%
\begin{equation}
(\cos \delta )\sum_{-\infty }^{n=-1}\big( \left\vert \Delta \chi
_{n}\right\vert ^{2}+q_{n}\left\vert \chi _{n}\right\vert ^{2}\big) \leq -%
{\rm Re}\big(\widehat{u}e^{-i\delta }\big).  \label{5.39}
\end{equation}%
It also follows that%
\begin{equation}
{\rm Re}\big(\widehat{u}e^{-i\delta }\big)<0.  \label{5.40}
\end{equation}%
Since $q_{n}\geq c>0,$ (\ref{5.39}) implies that for the solution $\chi
=(\chi _{n})$ def\/ined by (\ref{5.38}) we have (\ref{5.4}). Thus, the
statement of the theorem concerning the solution $\chi =(\chi _{n})$ is also
proved.

Finally, let us show that the solutions $\psi =(\psi _{n})$ and $\chi =(\chi
_{n})$ def\/ined by (\ref{5.24}) and (\ref{5.38}), respectively, are linearly
independent. We have%
\begin{equation}
W_{n}(\psi ,\chi )=W_{n}(\varphi +\widehat{v}\theta ,\varphi +\widehat{u}%
\theta )=(\widehat{u}-\widehat{v})W_{n}(\varphi ,\theta ).  \label{5.41}
\end{equation}%
Next, $W_{n}(\varphi ,\theta )\neq 0$ by (\ref{5.11}) and $\widehat{u}\neq
\widehat{v}$ by (\ref{5.26}), (\ref{5.40}). Therefore $W_{n}(\psi ,\chi
)\neq 0$ and hence $\psi $ and $\chi $ are linearly independent by Theorem %
\ref{Th4.5}.
\end{proof}


\section[The inverse operator $L^{-1}$]{The inverse operator $\boldsymbol{L^{-1}}$}\label{section6}

The following lemma will play crucial role in this and next sections.

\begin{lemma}
\label{Lem6.1} Let us set%
\begin{equation}
\sigma _{n}=\left\{
\begin{array}{lll}
e^{-i\delta } & \text{for} & n\leq -1, \\
e^{i\delta } & \text{for} & n\geq 0.%
\end{array}%
\right.  \label{6.1}
\end{equation}%
Under the conditions of Lemma {\rm \ref{Lem5.1}} the following formula holds:%
\begin{gather}
(\cos \delta )\left( \sum_{n=a}^{-1}+\sum_{n=2}^{b}\right) \big(
\left\vert \Delta y_{n}\right\vert ^{2}+q_{n}\left\vert y_{n}\right\vert
^{2}\big)  \nonumber \\
\qquad ={\rm Re}\left\{ \sigma _{n}(\Delta y_{n})\overline{y}_{n+1}\big|
_{a-1}^{b}+\left( \sum_{n=a}^{-1}+\sum_{n=2}^{b}\right) \sigma _{n}f_{n}%
\overline{y}_{n}\right\} .  \label{6.2}
\end{gather}
\end{lemma}

\begin{proof}
To prove (\ref{6.2}) we multiply (\ref{5.8}) by $e^{-i\delta }$ and (\ref%
{5.7}) by $e^{i\delta }$ and add together to get%
\begin{gather}
\left( e^{-i\delta }\sum_{n=a}^{-1}+e^{i\delta }\sum_{n=2}^{b}\right) \big(
\left\vert \Delta y_{n}\right\vert ^{2}+q_{n}\left\vert y_{n}\right\vert
^{2}\big)
\nonumber\\
\qquad =e^{-i\delta }(\Delta y_{n})\overline{y}_{n+1}\big|_{a-1}^{-1}+e^{i\delta
}(\Delta y_{n})\overline{y}_{n+1}\big|_{1}^{b}+\left( e^{-i\delta
}\sum_{n=a}^{-1}+e^{i\delta }\sum_{n=2}^{b}\right) f_{n}\overline{y}_{n}.
\label{6.3}
\end{gather}%
Next, using the conditions (\ref{5.6}) we have%
\begin{gather*}
e^{-i\delta }(\Delta y_{-1})\overline{y}_{0}-e^{i\delta }(\Delta y_{1})%
\overline{y}_{2}=e^{i\delta }(\Delta y_{1})\overline{y}_{0}-e^{i\delta
}(\Delta y_{1})\overline{y}_{2}\\
\qquad{}
=e^{i\delta }(\Delta y_{1})(\overline{y}_{0}-\overline{y}_{2})=e^{i\delta
}(\Delta y_{1})(\overline{y}_{0}-\overline{y}_{1}+\overline{y}_{1}-\overline{%
y}_{2})
\\
\qquad{}=e^{i\delta }(\Delta y_{1})(\overline{y}_{0}-\overline{y}_{-1}+\overline{y}%
_{1}-\overline{y}_{2})=e^{i\delta }(\Delta y_{1})(\overline{\Delta y_{-1}}-%
\overline{\Delta y_{1}})
\\
\qquad{}
=-e^{i\delta }\left\vert \Delta y_{1}\right\vert ^{2}+e^{i\delta }(\Delta
y_{1})e^{-2i\delta }\overline{\Delta y_{1}}=-e^{i\delta }\left\vert \Delta
y_{1}\right\vert ^{2}+e^{-i\delta }\left\vert \Delta y_{1}\right\vert ^{2}
=-2i(\sin \delta )\left\vert \Delta y_{1}\right\vert ^{2}.
\end{gather*}
Therefore taking (\ref{6.1}) into account we can rewrite (\ref{6.3}) in the
form%
\begin{gather}
\left( \sum_{n=a}^{-1}+\sum_{n=2}^{b}\right) \sigma _{n}\big( \left\vert
\Delta y_{n}\right\vert ^{2}+q_{n}\left\vert y_{n}\right\vert ^{2}\big)
\nonumber\\
\qquad{} =-2i(\sin \delta )\left\vert \Delta y_{1}\right\vert ^{2}+\sigma _{n}(\Delta
y_{n})\overline{y}_{n+1}\big|_{a-1}^{b}+\left(
\sum_{n=a}^{-1}+\sum_{n=2}^{b}\right) \sigma _{n}f_{n}\overline{y}_{n}.
\label{6.4}
\end{gather}%
Taking in (\ref{6.4}) the real parts of both sides and taking into account
that ${\rm Re}\,\sigma _{n}=\cos \delta $ for all $n$, we obtain (\ref{6.2}).
\end{proof}

Let $L:D\subset l_{0}^{2}\rightarrow l_{0}^{2}$ be the operator def\/ined
above in Section~\ref{section2}. Further, let $\psi =(\psi _{n})$ and $\chi
=(\chi _{n}),$ where $n\in
\mathbb{Z}
$, be solutions of problem (\ref{5.1}), (\ref{5.2}), constructed in Theorem~\ref{Th5.2}. Let us introduce the discrete Green function%
\begin{equation*}
G_{nk}=\frac{1}{W_{k}(\psi ,\chi )}\left\{
\begin{array}{lll}
\chi _{k}\psi _{n} & \text{if} & k\leq n, \\
\chi _{n}\psi _{k} & \text{if} & k\geq n,%
\end{array}%
\right. % \label{6.5}
\end{equation*}%
of discrete variables $k,n\in
\mathbb{Z}
.$ Note that by (\ref{5.41}) and (\ref{5.11}), we have%
\begin{equation}
W_{0}(\psi ,\chi )=-(\widehat{u}-\widehat{v})e^{2i\delta },\qquad
W_{k}(\psi ,\chi )=\left\{
\begin{array}{lll}
(\widehat{u}-\widehat{v})e^{2i\delta } & \text{if} & k\leq -1, \\
\widehat{u}-\widehat{v} & \text{if} & k\geq 1,%
\end{array}%
\right.  \label{6.6}
\end{equation}%
and, besides,%
\begin{equation*}
\widehat{u}\neq \widehat{v}  %\label{6.7}
\end{equation*}%
by (\ref{5.26}) and (\ref{5.40}).

\begin{theorem}
\label{Th6.1}Under the condition \eqref{5.3} the inverse operator $L^{-1}$
exists and is a bounded operator defined on the whole space $l_{0}^{2}.$
Next, for every $f=(f_{n})\in l_{0}^{2}$%
\begin{equation}
\big(L^{-1}f\big)_{n}=\sum_{k\in
\mathbb{Z}
_{0}}G_{nk}f_{k},\qquad n\in \mathbb{Z}
_{0},  \label{6.8}
\end{equation}%
and%
\begin{equation}
\big\Vert L^{-1}f\big\Vert \leq \frac{1}{c\cos \delta }\left\Vert
f\right\Vert \qquad \text{for all} \ \ f\in l_{0}^{2},  \label{6.9}
\end{equation}%
where $c$ is a constant from condition \eqref{5.3} and $\delta $ is from \eqref{5.2}, $\left\Vert \cdot \right\Vert $ denotes the norm of space $l_{0}^{2}.$
\end{theorem}

\begin{proof}
Let us show that
\[
\ker L=\left\{ y\in D:Ly=0\right\}
\]
consists only of the zero element. Indeed, if $y\in D$ and $Ly=0,$ then $%
(y)_{n\in
\mathbb{Z}
_{0}}$ satisf\/ies the equation%
\begin{equation}
-\Delta ^{2}y_{n-1}+q_{n}y_{n}=0,\qquad n\in
\mathbb{Z}
_{0},  \label{6.10}
\end{equation}%
in which $y_{0}$ and $y_{1}$ (these values arise in (\ref{6.10}) for $n=-1$
and $n=2,$ respectively) are def\/ined from the equations%
\begin{equation}
y_{-1}=y_{1},\qquad \Delta y_{-1}=e^{2i\delta }\Delta y_{1}.
\label{6.11}
\end{equation}%
Since $\chi $ and $\psi $ form a fundamental system of solutions of (\ref%
{6.10}), (\ref{6.11}), we can write
\[
y_{n}=C_{1}\psi _{n}+C_{2}\chi _{n},\qquad n\in
\mathbb{Z}
,
\]%
with some constants $C_{1}$ and $C_{2}.$ Hence%
\begin{equation}
W_{n}(y,\psi )=C_{1}W_{n}(\psi ,\psi )+C_{2}W_{n}(\chi ,\psi ),\qquad
n\in
\mathbb{Z}
.  \label{6.12}
\end{equation}%
Next, since $y\in l_{0}^{2},$ we have $y_{n}\rightarrow 0$ as $\left\vert
n\right\vert \rightarrow \infty $ and by (\ref{5.4}) we have $\psi
_{n}\rightarrow 0$ as $n\rightarrow \infty $ \ Hence $W_{n}(y,\psi
)\rightarrow 0$ as $n\rightarrow \infty .$ Besides $W_{n}(\psi ,\psi )=0$
for all $n$ and $W_{n}(\chi ,\psi )$ is equal to a nonzero constant for $%
n\geq 1$ by (\ref{6.6}). Therefore taking the limit in (\ref{6.12}) as $%
n\rightarrow \infty $ we get that $C_{2}=0.$ It can similarly be shown, by
considering $W_{n}(y,\chi )$, that $C_{1}=0.$ Thus $y=0.$

It follows that the inverse operator $L^{-1}$ exists. Now take an arbitrary $%
f=(f_{n})_{n\in
\mathbb{Z}
_{0}}\in l_{0}^{2}$ and extend the sequence $(f_{n})_{n\in
\mathbb{Z}
_{0}}$ to the values $n=0$ and $n=1$ by setting%
\begin{equation*}
f_{0}=f_{1}=0.  %\label{6.13}
\end{equation*}%
Let us put%
\begin{gather}
g_{n}=\sum_{k\in
\mathbb{Z}
_{0}}G_{nk}f_{k}=\sum_{k\in
\mathbb{Z}
}G_{nk}f_{k}
=\psi _{n}\sum_{-\infty }^{k=n}\frac{\chi _{k}f_{k}}{W_{k}(\psi ,\chi )}%
+\chi _{n}\sum_{k=n+1}^{\infty }\frac{\psi _{k}f_{k}}{W_{k}(\psi ,\chi )},\qquad
n\in
\mathbb{Z}
.  \label{6.14}
\end{gather}%
Then it is easy to check that this sequence $(g_{n}),$ where $n\in
\mathbb{Z}
$, satisf\/ies the equations%
\begin{gather}
-\Delta ^{2}g_{n-1}+q_{n}g_{n}=f_{n},\qquad n\in
\mathbb{Z}
_{0},  \label{6.15}
\\
g_{-1}=g_{1},\qquad \Delta g_{-1}=e^{2i\delta }\Delta g_{1}.
\label{6.16}
\end{gather}%
We want to show that $g=(g_{n})_{n\in
\mathbb{Z}
_{0}}\in l_{0}^{2}$ and that%
\begin{equation}
\left\Vert g\right\Vert \leq \frac{1}{c\cos \delta }\left\Vert f\right\Vert .
\label{6.17}
\end{equation}%
For this purpose we take the sequences of integers $a_{m}$ and $b_{m}$
def\/ined for any positive integer~$m$, such that%
\[
a_{m}<0<b_{m}\qquad \text{and} \qquad a_{m}\rightarrow -\infty ,\qquad
b_{m}\rightarrow \infty \qquad \text{as} \qquad m\rightarrow \infty .
\]%
Next, for each $m$ we def\/ine the sequence $(f_{n}^{(m)})_{n\in
\mathbb{Z}
}$ by%
\begin{gather}
f_{n}^{(m)}=f_{n}\qquad \text{if} \qquad a_{m}\leq n\leq b_{m},  \label{6.18}
\\
f_{n}^{(m)}=0 \qquad \text{if} \qquad n<a_{m} \qquad \text{or} \qquad n>b_{m},  \label{6.19}
\end{gather}%
and put%
\begin{gather*}
g_{n}^{(m)}=\sum_{k\in
\mathbb{Z}
_{0}}G_{nk}f_{k}^{(m)}=\sum_{k\in
\mathbb{Z}
}G_{nk}f_{k}^{(m)}
=\psi _{n}\sum_{-\infty }^{k=n}\frac{\chi _{k}f_{k}^{(m)}}{W_{k}(\psi ,\chi )%
}+\chi _{n}\sum_{k=n+1}^{\infty }\frac{\psi _{k}f_{k}^{(m)}}{W_{k}(\psi
,\chi )},\qquad n\in \mathbb{Z}
.  %\label{6.20}
\end{gather*}%
It follows that%
\begin{gather}
g_{n}^{(m)}=\left\{
\begin{array}{lll}
\displaystyle \chi _{n}\sum_{k=a_{m}}^{b_{m}}\frac{\psi _{k}f_{k}}{W_{k}(\psi ,\chi )} &
\text{if} & n<a_{m}, \vspace{1mm}\\
\displaystyle \psi _{n}\sum_{k=a_{m}}^{b_{m}}\frac{\chi _{k}f_{k}}{W_{k}(\psi ,\chi )} &
\text{if} & n>b_{m}.%
\end{array}%
\right.  \label{6.21}
\end{gather}%
We have also that, for each $m,$%
\begin{gather}
-\Delta ^{2}g_{n-1}^{(m)}+q_{n}g_{n}^{(m)}=f_{n}^{(m)},\qquad n\in
\mathbb{Z}
_{0},  \label{6.22}
\\
g_{-1}^{(m)}=g_{1}^{(m)},\qquad \Delta g_{-1}^{(m)}=e^{2i\delta }\Delta
g_{1}^{(m)}.  \label{6.23}
\end{gather}%
Fix $m$ and take a positive integer $N$ such that%
\[
-N<a_{m}\qquad \text{and} \qquad b_{m}<N.
\]%
Then applying Lemma \ref{Lem6.1} to (\ref{6.22}), (\ref{6.23}) we can write%
\begin{gather}
(\cos \delta )\left( \sum_{n=-N}^{-1}+\sum_{n=2}^{N}\right) \big(
\big\vert \Delta g_{n}^{(m)}\big\vert ^{2}+q_{n}\big\vert
g_{n}^{(m)}\big\vert ^{2}\big)  \nonumber \\
\qquad {}= {\rm Re}\left\{ \sigma _{n}(\Delta g_{n}^{(m)})\overline{g}%
_{n+1}^{(m)}\big| _{-N-1}^{N}+\left( \sum_{n=-N}^{-1}+\sum_{n=2}^{N}\right)
\sigma _{n}f_{n}^{(m)}\overline{g}_{n}^{(m)}\right\} .  \label{6.24}
\end{gather}%
It follows from (\ref{6.21}) by (\ref{5.4}) that%
\[
\sum_{n\in
\mathbb{Z}
}\big\vert g_{n}^{(m)}\big\vert ^{2}<\infty.
\]
Therefore the sums on the right-hand side of (\ref{6.24}) are convergent as $%
N\rightarrow \infty $ and besides%
\[
{\rm Re}\big\{ \sigma _{n}(\Delta g_{n}^{(m)})\overline{g}_{n+1}^{(m)}\big|
_{-N-1}^{N}\big\} \rightarrow 0\qquad \text{as} \ \  N\rightarrow \infty .
\]%
(note that $\left\vert \sigma _{n}\right\vert =1$ for all $n$ by (\ref{6.1}%
)). Now taking the limit in (\ref{6.24}) as $N\rightarrow \infty ,$ we get%
\begin{equation}
(\cos \delta )\sum_{n\in
\mathbb{Z}
_{0}}\big( \big\vert \Delta g_{n}^{(m)}\big\vert ^{2}+q_{n}\big\vert
g_{n}^{(m)}\big\vert ^{2}\big) ={\rm Re}\sum_{n\in
\mathbb{Z}
_{0}}\sigma _{n}f_{n}^{(m)}\overline{g}_{n}^{(m)}.  \label{6.25}
\end{equation}%
Using the condition (\ref{5.3}) we get from (\ref{6.25}) that%
\begin{gather*}
\sum_{n\in
\mathbb{Z}
_{0}}\big\vert g_{n}^{(m)}\big\vert ^{2}\leq \frac{1}{c\cos \delta }%
{\rm Re}\sum_{n\in
\mathbb{Z}
_{0}}\sigma _{n}f_{n}^{(m)}\overline{g}_{n}^{(m)}
\leq \frac{1}{c\cos \delta }\left\vert \sum_{n\in
\mathbb{Z}
_{0}}\sigma _{n}f_{n}^{(m)}\overline{g}_{n}^{(m)}\right\vert \\
\qquad {} \leq \frac{1}{%
c\cos \delta }\sum_{n\in
\mathbb{Z}
_{0}}\big\vert f_{n}^{(m)}\overline{g}_{n}^{(m)}\big\vert
\leq \frac{1}{c\cos \delta }\left\{ \sum_{n\in
\mathbb{Z}
_{0}}\big\vert f_{n}^{(m)}\big\vert ^{2}\right\} ^{1/2}\left\{ \sum_{n\in
\mathbb{Z}
_{0}}\big\vert g_{n}^{(m)}\big\vert ^{2}\right\} ^{1/2}.
\end{gather*}
Hence%
\[
\left\{ \sum_{n\in
\mathbb{Z}
_{0}}\big\vert g_{n}^{(m)}\big\vert ^{2}\right\} ^{1/2}\leq \frac{1}{%
c\cos \delta }\left\{ \sum_{n\in
\mathbb{Z}
_{0}}\big\vert f_{n}^{(m)}\big\vert ^{2}\right\} ^{1/2},
\]%
that is,%
\begin{equation}
\big\Vert g^{(m)}\big\Vert \leq \frac{1}{c\cos \delta }\big\Vert
f^{(m)}\big\Vert .  \label{6.26}
\end{equation}

Writing (\ref{6.22}), (\ref{6.23}) for $m=m_{1}$ and $m=m_{2},$ and
subtracting the obtained equations side-by-side, we get%
\begin{gather*}
-\Delta
^{2}\big[g_{n-1}^{(m_{1})}-g_{n-1}^{(m_{2})}\big]+q_{n}\big[g_{n}^{(m_{1})}-g_{n}^{(m_{2})}\big]=f_{n}^{(m_{1})}-f_{n}^{(m_{2})},%
\qquad n\in
\mathbb{Z}
_{0},
\\
g_{-1}^{(m_{1})}-g_{-1}^{(m_{2})}=g_{1}^{(m_{1})}-g_{1}^{(m_{2})},\qquad \Delta \big[ g_{-1}^{(m_{1})}-g_{-1}^{(m_{2})}\big]=e^{2i\delta }\Delta
\big[ g_{1}^{(m_{1})}-g_{1}^{(m_{2})}\big].
\end{gather*}
Hence, repeating the same reasonings as above, we get%
\[
\big\Vert g^{(m_{1})}-g^{(m_{2})}\big\Vert \leq \frac{1}{c\cos \delta }%
\big\Vert f^{(m_{1})}-f^{(m_{2})}\big\Vert .
\]%
It follows that $g^{(m)}$ converges in $l_{0}^{2}$ to an element $\widetilde{%
g}$ as $m\rightarrow \infty .$ On the other hand, it can be seen from (\ref{6.14}), (\ref{6.21}) taking into account (\ref{6.18}), (\ref{6.19}) that%
\[
g_{n}^{(m)}\rightarrow g_{n}\qquad \text{as} \ \ m\rightarrow \infty ,
\]%
for each $n$. Consequently, $\widetilde{g}=g$ and hence $g\in l_{0}^{2}.$
Passing in (\ref{6.26}) to the limit as $m\rightarrow \infty ,$ we get (\ref{6.17}).

Next, from (\ref{6.15}) we have%
\[
q_{n}g_{n}=f_{n}+g_{n-1}-2g_{n}+g_{n+1},\qquad n\in
\mathbb{Z}
_{0}.
\]
Hence $(q_{n}g_{n})_{n\in
\mathbb{Z}
_{0}}\in l_{0}^{2}.$ Therefore $g\in D,$ where $D$ is the domain of the
operator~$L$. If we def\/ine an operator $B:l_{0}^{2}\rightarrow l_{0}^{2}$ by
the formula $Bf=g,$ where $f=(f_{n})_{n\in
\mathbb{Z}
_{0}}\in l_{0}^{2}$ and $g=(g_{n})_{n\in
\mathbb{Z}
_{0}}$ with~$g_{n}$ def\/ined by (\ref{6.14}), then we get by (\ref{6.15}), (\ref{6.16}) that $LBf=f.$ Therefore $B$ is the inverse of the operator $%
L:B=L^{-1},$ so that $g=L^{-1}f$ and from (\ref{6.14}) and (\ref{6.17}) we
get~(\ref{6.8}) and~(\ref{6.9}), respectively.
\end{proof}

\section[Completely continuity of the operator $L^{-1}$]{Completely continuity of the operator $\boldsymbol{L^{-1}}$}\label{section7}


In this section we will show that the operator $L^{-1}$ is completely
continuous, that is, it is continuous and maps bounded sets into relatively
compact sets.

\begin{theorem}
\label{Th7.1}Let%
\begin{equation}
q_{n}\geq c>0\qquad \text{for} \ \ n\in
\mathbb{Z}_{0},  \label{7.1}
\end{equation}%
and
\begin{equation}
\lim_{\left\vert n\right\vert \rightarrow \infty }q_{n}=\infty .  \label{7.2}
\end{equation}%
Then the operator $L^{-1}$ is completely continuous.
\end{theorem}

\begin{proof}
The operator $L^{-1}$ is continuous in virtue of (\ref{6.9}) that holds
under the condition (\ref{7.1}). In order to show that $L^{-1}$ maps bounded
sets into relatively compact sets consider any bounded set $X$ in $l_{0}^{2}$,
\[
X=\big\{ f\in l_{0}^{2}:\left\Vert f\right\Vert \leq d\big\} ,
\]%
and prove that $L^{-1}(X)=Y$ is relatively compact in $l_{0}^{2}.$ To this
end, we use the following known (see~\cite{[17]}) criterion for the relative
compactness in $l_{0}^{2}:$ \textit{A set $Y\subset l_{0}^{2}$
is relatively compact if and only if  $Y$ is
bounded and for every $\varepsilon >0$ there exists a positive
integer $n_{0}$ (depending only on~$\varepsilon $) such
that}
\[
\sum_{\left\vert n\right\vert >n_{0}}\left\vert y_{n}\right\vert ^{2}\leq
\varepsilon \qquad \text{\textit{for all}} \ \ y\in Y.
\]

Take an arbitrary $f\in X$ and set
\[
L^{-1}f=y.
\]
Then $Ly=f$ or explicitly%
\begin{equation}
-\Delta ^{2}y_{n-1}+q_{n}y_{n}=f_{n},\qquad n\in
\mathbb{Z}
_{0},  \label{7.3}
\end{equation}%
where $y_{0}$ and $y_{1}$ are def\/ined from the equations%
\begin{equation}
y_{-1}=y_{1},\qquad \Delta y_{-1}=e^{2i\delta }\Delta y_{1}.
\label{7.4}
\end{equation}
Note that $y_{0}$ and $y_{1}$ are needed when we write out equation (\ref{7.3}) for $n=-1$ and $n=2,$ respectively.

Applying Lemma \ref{Lem6.1} to (\ref{7.3}), (\ref{7.4}), we get that for any
integers $a\leq -1$ and $b\geq 2,$%
\begin{gather}
(\cos \delta )\left( \sum_{n=a}^{-1}+\sum_{n=2}^{b}\right) \big(
\left\vert \Delta y_{n}\right\vert ^{2}+q_{n}\left\vert y_{n}\right\vert
^{2}\big)  \nonumber \\
\qquad {} ={\rm Re}\left\{ \sigma _{n}(\Delta y_{n})\overline{y}_{n+1}\big|_{a-1}^{b}+\left( \sum_{n=a}^{-1}+\sum_{n=2}^{b}\right) \sigma _{n}f_{n}%
\overline{y}_{n}\right\} ,  \label{7.5}
\end{gather}
where $\sigma _{n}$ is def\/ined by (\ref{6.1}).

Since $f,y\in l_{0}^{2}$ and $\left\vert \sigma _{n}\right\vert =1,$ the
sums on the right-hand side of (\ref{7.5}) are convergent as $a\rightarrow
-\infty ,$ $b\rightarrow \infty .$ Also from $y\in l_{0}^{2}$ it follows
that $y_{n}\rightarrow 0$ as $\left\vert n\right\vert \rightarrow \infty $
so that%
\[
\sigma _{n}(\Delta y_{n})\overline{y}_{n+1}\big|_{a-1}^{b}\rightarrow 0\qquad \text{as} \qquad a\rightarrow -\infty ,\qquad b\rightarrow \infty .
\]
Consequently, we arrive at the equality%
\[
(\cos \delta )\sum_{n\in
\mathbb{Z}
_{0}}\big( \left\vert \Delta y_{n}\right\vert ^{2}+q_{n}\left\vert
y_{n}\right\vert ^{2}\big) ={\rm Re}\sum_{n\in
\mathbb{Z}
_{0}}\sigma _{n}f_{n}\overline{y}_{n}.
\]%
Hence%
\begin{equation}
(\cos \delta )\sum_{n\in
\mathbb{Z}
_{0}}q_{n}\left\vert y_{n}\right\vert ^{2}\leq {\rm Re}\sum_{n\in
\mathbb{Z}
_{0}}\sigma _{n}f_{n}\overline{y}_{n}  \label{7.6}
\end{equation}%
and therefore using (\ref{7.1}) and $\left\Vert f\right\Vert \leq d,$ we get{\samepage
\begin{equation}
\left\Vert y\right\Vert \leq \frac{d}{c\cos \delta }\qquad \text{for all} \ \
y\in Y.  \label{7.7}
\end{equation}%
This means that the set $Y=L^{-1}(X)$ is bounded.}

From (\ref{7.6}) we also have, using (\ref{7.7}),%
\begin{equation}
\sum_{n\in \mathbf{%
%TCIMACRO{\U{2124} }%
%BeginExpansion
\mathbb{Z}
%EndExpansion
}_{0}}q_{n}\left\vert y_{n}\right\vert ^{2}\leq \frac{d^{2}}{c\cos
^{2}\delta } \qquad \text{for all} \ \ y\in Y.  \label{7.8}
\end{equation}

Take now an arbitrary $\varepsilon >0$. By condition (\ref{7.2}) we can
choose a positive integer $n_{0}$ such that
\[
q_{n}\geq \frac{d^{2}}{\varepsilon c\cos ^{2}\delta } \qquad \text{for} \ \
\left\vert n\right\vert >n_{0}.
\]%
Then we get from (\ref{7.8}) that
\[
\sum_{\left\vert n\right\vert >n_{0}}\left\vert y_{n}\right\vert ^{2}\leq
\varepsilon \qquad \text{for all} \ \ y\in Y.
\]

Thus the completely continuity of the operator $L^{-1}$ is proved.
\end{proof}
